\documentclass[10pt,a4paper]{amsart}
\usepackage[margin=19mm]{geometry}
\usepackage{amsmath,amssymb,mathtools}
\usepackage[scr=rsfs,cal=euler]{mathalfa}
\usepackage{xfrac}
\usepackage[unicode,hidelinks,bookmarks=false]{hyperref}
\newcommand{\ie}{i.e.\ }
\newcommand{\cf}{cf.\ }
\newcommand{\resp}{respectively\ }
\newcommand{\Subsection}{\subsection}
\providecommand{\nobreakdash}{\nobreak\hbox{-}}
\setlength{\emergencystretch}{2em}
\multlinegap0pt
\usepackage{amssymb}
\usepackage{float}
\usepackage{xcolor}
\usepackage[shortlabels]{enumitem}
\usepackage{algorithm}\usepackage{algpseudocode}
\newtheorem{theorem}{Theorem}
\newtheorem{remark}{Remark}
\title{Well-posedness and asymptotic behavior of solutions to a second order nonlocal parabolic MEMS equation}
\author{Yufei Wei, Yanyan Zhang}
\date{Source: arXiv:2603.07013v1 (CC BY 4.0)}
\begin{document}
\maketitle

\noindent\emph{Source and licence.} Well-posedness and asymptotic behavior of solutions to a second order nonlocal parabolic MEMS equation. Version arXiv:2603.07013v1 (CC BY 4.0), distributed by arXiv under the Creative Commons Attribution 4.0 International licence, http://creativecommons.org/licenses/by/4.0/, as recorded in the arXiv licence metadata for this exact version and shown on the abstract page https://arxiv.org/abs/2603.07013v1. Full licence notice: \texttt{licenses/arxiv-2603.07013-CC-BY-4.0.txt}.\par\smallskip

\noindent\textit{Editorial scope.} Counted unit: the article's theorem theorem1.2 (source0.tex lines 271--282). The article's complete original proof of it is delivered with it (source0.tex lines 1358--1531, one \texttt{proof} environment of 174 lines, which the article places away from the statement). No mathematics is written, repaired or re-derived: the statement and the proof are the article's own lines, in the article's own notation, and no retained line is substituted or re-flowed.  Retained uncounted as the source-local context the proof needs: lines 144--153; lines 254--256; lines 857--859; lines 1066--1073; lines 1070--1072; lines 1241--1247; lines 1283--1314.  Nothing was omitted from any retained span, so the package carries no omission record.  No retained line contains a citation, so the wrapper carries no bibliography and no bibliography entry.  No retained line contains a figure environment, an image-inclusion command or drawing code, so no image is delivered and the package declares no asset dependency.  Editorial formatting only, in addition: an AMS \texttt{amsart} preamble that restates the article's own declarations and macros used by the retained lines, namely the environments \texttt{theorem}, \texttt{remark} on separate counters, together with the standard packages those lines need.  The retained environments print in this document as Theorem 1, Remark 1 in order of appearance, whereas the article prints the same items with the numbering of its own full document.  The complete native source of the article as distributed in the arXiv e-print for this exact version is bundled unchanged as \texttt{sources/195911/source0.tex}.
To simplify the problem, we introduce dimensionless variables and adjust parameters to normalize the elastic coefficient $T$, damping coefficient $a$, and gap distance $d$. By neglecting membrane rigidity and thickness (i.e. $A = B = 0$), we obtain the equation:
\begin{equation}
    \noindent
    \begin{cases}
        u_{t}-\Delta u=\frac{\lambda}{(1-u)^2(1+\alpha\int_{\Omega}\frac{1}{1-u}\,\mathrm{d}x)^2},&x\in\Omega,t>0\\
        u=0,&x\in\partial\Omega,t>0\\
        u(x,0)=u_0(x),&x\in\Omega
    \end{cases}\label{1.1}
\end{equation}\\
where $\lambda,\alpha>0$,\ $\Omega\subset\mathbb{R}^{N}$. Without loss of generality, we assume $\alpha = 1$. Equation (\ref{1.1}) is the second order nonlocal parabolic MEMS equation studied in this paper.

  \begin{equation}
      \lim_{t \to +\infty} \|u(\cdot,t)-\psi\|_{H^2}=0. \label{1.6}
  \end{equation}

\begin{equation}
  \frac{\,\mathrm{d}E}{\,\mathrm{d}t}=-\int_{\Omega}u_{t}^2\,\mathrm{d}x\le 0.\label{2.2}
\end{equation}

\begin{theorem}\label{theorem2.1}
    Let $\mathscr{E}$ be the set of equilibrium points of equation (\ref{1.1}):
    $$\mathscr{E}=\left\{\phi\mid -\Delta\phi=\frac{\lambda}{(1-\phi)^2(1+\int_{\Omega}\frac{1}{1-\phi}\,\mathrm{d}x)^2},\  \phi|_{\partial\Omega}=0\right\},$$
    and let $\psi\in\mathscr{E}$. Then there exist $\sigma>0$ depending on $\psi$ and $\theta\in (0,\frac{1}{2})$ such that for every $u\in H^2\cap H_0^1(\Omega)$ satisfying $\|u-\psi\|_{H^2}\le\sigma$:
    \begin{equation}
        \left\|\Delta u+\frac{\lambda}{(1-u)^2(1+\int_{\Omega}\frac{1}{1-u}\,\mathrm{d}x)^2}\right\|_{L^2}\ge |E(u)-E(\psi)|^{1-\theta}.\label{2.11}
    \end{equation}
\end{theorem}

    \begin{equation}
        \left\|\Delta u+\frac{\lambda}{(1-u)^2(1+\int_{\Omega}\frac{1}{1-u}\,\mathrm{d}x)^2}\right\|_{L^2}\ge |E(u)-E(\psi)|^{1-\theta}.
    \end{equation}

\begin{remark}\label{remark2.1}
    If we modify (\ref{2.11}) to
    \begin{equation}
        \left\|\Delta u+\frac{\lambda}{(1-u)^2(1+\int_{\Omega}\frac{1}{1-u}\,\mathrm{d}x)^2}\right\|_{L^{2}}\ge C|E(u)-E(\psi)|^{1-\theta}\label{2.25}
    \end{equation}
    where $C>0$ is a constant, then $\theta$ can be taken as $\frac{1}{2}$.
\end{remark}

    By Theorem \ref{theorem2.1}, for $\psi\in\omega(u_0)$, there exist $\sigma>0$ and $\theta\in (0,\frac{1}{2})$ depending on $\psi$ such that when $\|u-\psi\|_{H^2}\le\sigma$:
    \begin{equation}
        \left\|\Delta u+\frac{\lambda}{(1-u)^2(1+\int_{\Omega}\frac{1}{1-u}\,\mathrm{d}x)^2}\right\|_{L^{2}}\ge |E(u)-E(\psi)|^{1-\theta}.\label{2.26}
    \end{equation}
    Note that (\ref{2.2}) is equivalent to
    \begin{equation}
        \frac{\mathrm{d}}{\mathrm{d}t}(E(u(t))-E(\psi))+\|u_t\|_{L^2}\left\|-\Delta u-\frac{\lambda}{(1-u)^2(1+\int_{\Omega}\frac{1}{1-u}\,\mathrm{d}x)^2}\right\|_{L^2}=0,\label{2.27}
    \end{equation}
    and
    \begin{equation}
        \frac{\mathrm{d}}{\mathrm{d}t}(E(u(t))-E(\psi))^{\theta}=\theta(E(u(t))-E(\psi))^{\theta-1} \frac{\mathrm{d}}{\mathrm{d}t}(E(u(t))-E(\psi)).\label{2.28}
    \end{equation}
    Substituting (\ref{2.27}) into (\ref{2.28}) yields:
    \begin{equation}
        \begin{aligned}
            &\quad\ \ \frac{\mathrm{d}}{\mathrm{d}t}(E(u(t))-E(\psi))^{\theta}\\
            &=\ -\theta(E(u(t))-E(\psi))^{\theta-1}\|u_t\|_{L^2}\left\|-\Delta u-\frac{\lambda}{(1-u)^2(1+\int_{\Omega}\frac{1}{1-u}\,\mathrm{d}x)^2}\right\|_{L^2}.
        \end{aligned}\label{2.29}
    \end{equation}
    If there exists a sufficiently large $T>0$ such that $\forall\ t>T$, $\|u-\psi\|_{H^2}<\sigma$, then by (\ref{2.26}):
    \begin{align*}
        &\quad\ \ -\frac{1}{\theta}\frac{\mathrm{d}}{\mathrm{d}t}(E(u(t))-E(\psi))^{\theta}\\
        &=\ (E(u(t))-E(\psi))^{\theta-1}\|u_t\|_{L^2}\left\|-\Delta u-\frac{\lambda}{(1-u)^2(1+\int_{\Omega}\frac{1}{1-u}\,\mathrm{d}x)^2}\right\|_{L^2}\\
        &\ge\ \|u_t\|_{L^2}.
    \end{align*}
    Therefore,
    $$\int_{T}^{+\infty}\|u_t\|_{L^2}\,\mathrm{d}\tau <+\infty.$$
    Moreover, since
    $$\|u(t)-u(s)\|_{L^2}=\left\|\int_{s}^{t}u_t\,\mathrm{d}\tau\right\|_{L^2}\le \int_{s}^{t}\|u_t\|_{L^2}\,\mathrm{d}\tau,$$
    we have
    $$\|u(t)-\psi\|_{L^2}\to 0,\quad t\to +\infty.$$
    By the relative compactness of $\bigcup_{t\ge T}u(t)$ in $H^2\cap H_0^1(\Omega)$, (\ref{1.6}) is proved.

\begin{theorem}\label{theorem1.2}
    Let $\theta\in (0,\frac{1}{2}]$ be defined as in Theorem \ref{theorem2.1}. If $\theta=\frac{1}{2}$, then the solution $u$ of equation (\ref{1.1}) decays exponentially to $\psi$ in $H^2\cap H_0^1(\Omega)$, i.e.
    \begin{equation}
        \|u(\cdot,t)-\psi\|_{H^2}\le C_0e^{-C_1t},\quad t\to +\infty,\label{2.33}
    \end{equation}
    where $C_0,\ C_1>0$ are constants.
    If $\theta\in (0,\frac{1}{2})$, then $u$ decays polynomially to $\psi$ in $H^2\cap H_0^1(\Omega)$, i.e.
    \begin{equation}
        \|u(\cdot,t)-\psi\|_{H^2}\le C(1+t)^{-\frac{\theta}{1-2\theta}},\quad t\to +\infty,\label{2.34}
    \end{equation}
    where $C>0$ is a constant.
\end{theorem}

\begin{proof}[Proof of Theorem \ref{theorem1.2}]
First, we prove that the solution $u$ of equation (\ref{1.1}) satisfies in $L^2(\Omega)$:
\begin{itemize}
    \item When $\theta=\frac{1}{2}$: $\|u(\cdot,t)-\psi\|_{L^2}\le C_0e^{-C_1t}$ as $t\to +\infty$.
    \item When $\theta\in (0,\frac{1}{2})$: $\|u(\cdot,t)-\psi\|_{L^2}\le C(1+t)^{-\frac{\theta}{1-2\theta}}$ as $t\to +\infty$.
\end{itemize}
  Indeed, by Theorem \ref{theorem2.1} and Remark \ref{remark2.1}, when $t$ is sufficiently large:
  \begin{equation}
    \|u_t\|_{L^2}= \left\|\Delta u+\frac{\lambda}{(1-u)^2(1+\int_{\Omega}\frac{1}{1-u}\,\mathrm{d}x)^2}\right\|_{L^{2}}\ge C|E(u)-E(\psi)|^{1-\theta},\quad \theta\in\left (0,\frac{1}{2}\right ].\label{2.35}
  \end{equation}
  Let $y(t)=(E(u)-E(\psi))^{\theta}$. From (\ref{2.29}) and (\ref{2.35}):
  \begin{equation}
      \frac{\,\mathrm{d}y}{\,\mathrm{d}t}+C\|u_t\|_{L^2}\le 0.\label{2.36}
  \end{equation}
  Integrating (\ref{2.36}):
  $$\int_{t}^{+\infty}\|u_t\|_{L^2}\,\mathrm{d}\tau\le \frac{1}{C}y(t),$$
  thus
  \begin{equation}
      \|u(\cdot,t)-\psi\|_{L^2}\le\left\|\int_{t}^{+\infty}u_t\,\mathrm{d}\tau\right\|_{L^2}\le\int_{t}^{+\infty}\|u_t\|_{L^2}\,\mathrm{d}\tau\le \frac{1}{C}y(t).\label{2.37}
  \end{equation}
  Combining (\ref{2.35}) and (\ref{2.36}) yields the differential inequality:
  \begin{equation}
      \frac{\,\mathrm{d}y}{\,\mathrm{d}t}+Cy^{\frac{1-\theta}{\theta}}\le 0.\label{2.38}
  \end{equation}
  \noindent
  When $\theta=\frac{1}{2}$, exponential decay follows directly from (\ref{2.37}) and (\ref{2.38}):
  $$\|u(\cdot,t)-\psi\|_{L^2}\le C_0e^{-C_1t}.$$
  \noindent
  When $\theta\in (0,\frac{1}{2})$, consider the equation:
  \begin{equation}
      \begin{cases}
          \frac{\,\mathrm{d}z}{\,\mathrm{d}t}+Cz^{\frac{1-\theta}{\theta}}=0 \\
          z(0)=y(0)>0 \label{2.39}
      \end{cases}
  \end{equation}
  The solution to (\ref{2.39}) is:
  $$z(t)=\left(-C\frac{2\theta-1}{\theta}t+y(0)^{\frac{2\theta-1}{\theta}}\right)^{\frac{\theta}{2\theta-1}}.$$
  Let $w(t)=z(t)-y(t)$. Subtracting (\ref{2.39}) from (\ref{2.38}):
  $$\frac{\,\mathrm{d}w}{\,\mathrm{d}t}+C(z^{\frac{1-\theta}{\theta}}-y^{\frac{1-\theta}{\theta}})\ge 0.$$
  Applying the mean value theorem:
  \begin{equation}
      \begin{cases}
          \frac{\,\mathrm{d}w}{\,\mathrm{d}t}+C\frac{1-\theta}{\theta}\xi^{\frac{1-2\theta}{\theta}}w\ge 0 \\
          w(0)=0 \label{2.40}
      \end{cases}
  \end{equation}
  where $\xi$ lies between $y$ and $z$.\\
  Multiplying (\ref{2.40}) by $e^{\int_{0}^{t}C\frac{1-\theta}{\theta}\xi(\tau)^{\frac{1-2\theta}{\theta}}\,\mathrm{d}\tau}$:
  $$\frac{\mathrm{d}}{\mathrm{d}t}\left(w(t)e^{\int_{0}^{t}C\frac{1-\theta}{\theta}\xi(\tau)^{\frac{1-2\theta}{\theta}}\,\mathrm{d}\tau}\right)\ge 0,$$
  and since $w(0)=0$:
  $$w(t)e^{\int_{0}^{t}C\frac{1-\theta}{\theta}\xi(\tau)^{\frac{1-2\theta}{\theta}}\,\mathrm{d}\tau}\ge 0,$$
  thus
  $$w(t)\ge 0,$$
  and therefore
  $$\|u(\cdot,t)-\psi\|_{L^2}\le \frac{1}{C}y(t)\le \frac{1}{C}z(t)=C(1+t)^{-\frac{\theta}{1-2\theta}}.$$

  Next, we prove that in $H^2\cap H_0^1(\Omega)$, (\ref{2.33}) and (\ref{2.34}) hold for $\theta=\frac{1}{2}$ and $\theta\in (0,\frac{1}{2})$ respectively.
  Let $v=u-\psi$. Then $v$ satisfies:
  \begin{equation}
    \begin{cases}
        v_{t}-\Delta v=\frac{\lambda}{(1-u)^2(1+\int_{\Omega}\frac{1}{1-u}\,\mathrm{d}x)^2}-\frac{\lambda}{(1-\psi)^2(1+\int_{\Omega}\frac{1}{1-\psi}\,\mathrm{d}x)^2},&x\in\Omega,t>0\\
        v=0,&x\in\partial\Omega,t>0\\
        v(x,0)=u_0(x)-\psi(x),&x\in\Omega
    \end{cases}\label{4.1}
\end{equation}
  Multiplying (\ref{4.1}) by $v$ and integrating over $\Omega$:
\begin{align*}
    &\quad\ \ \frac{1}{2}\frac{\mathrm{d}}{\mathrm{d}t}\|v\|_{L^2}^2+\|\nabla v\|_{L^2}^2\\
    &\ =\int_{\Omega}\left(\frac{\lambda}{(1-u)^2(1+\int_{\Omega}\frac{1}{1-u}\,\mathrm{d}x)^2}-\frac{\lambda}{(1-\psi)^2(1+\int_{\Omega}\frac{1}{1-\psi}\,\mathrm{d}x)^2}\right)v\,\mathrm{d}x,
\end{align*}
where
\begin{align*}
    &\quad\frac{\lambda}{(1-u)^2(1+\int_{\Omega}\frac{1}{1-u}\,\mathrm{d}x)^2}-\frac{\lambda}{(1-\psi)^2(1+\int_{\Omega}\frac{1}{1-\psi}\,\mathrm{d}x)^2}\\
    &\ =\lambda\left[\frac{1}{(1-u)^2(1+\int_{\Omega}\frac{1}{1-u}\,\mathrm{d}x)^2}-\frac{1}{(1-u)^2(1+\int_{\Omega}\frac{1}{1-\psi}\,\mathrm{d}x)^2}\right.\\
    &\quad\ \left.+\frac{1}{(1-u)^2(1+\int_{\Omega}\frac{1}{1-\psi}\,\mathrm{d}x)^2}-\frac{1}{(1-\psi)^2(1+\int_{\Omega}\frac{1}{1-\psi}\,\mathrm{d}x)^2}\right]\\
    &\ =\lambda\left[\frac{1}{(1-u)^2}\frac{(1+\int_{\Omega}\frac{1}{1-\psi}\,\mathrm{d}x)^2-(1+\int_{\Omega}\frac{1}{1-u}\,\mathrm{d}x)^2}{(1+\int_{\Omega}\frac{1}{1-u}\,\mathrm{d}x)^2(1+\int_{\Omega}\frac{1}{1-\psi}\,\mathrm{d}x)^2}\right.\\
    &\quad\ \left.+\frac{1}{(1+\int_{\Omega}\frac{1}{1-\psi}\,\mathrm{d}x)^2}\frac{(1-\psi)^2-(1-u)^2}{(1-u)^2(1-\psi)^2}\right]\\
    &\ =\lambda\left[\frac{1}{(1-u)^2}\frac{(2+\int_{\Omega}\frac{1}{1-u}\,\mathrm{d}x+\int_{\Omega}\frac{1}{1-\psi}\,\mathrm{d}x)(\int_{\Omega}\frac{-v}{(1-u)(1-\psi)}\,\mathrm{d}x)}{(1+\int_{\Omega}\frac{1}{1-u}\,\mathrm{d}x)^2(1+\int_{\Omega}\frac{1}{1-\psi}\,\mathrm{d}x)^2}\right.\\
    &\quad\ \left.+\frac{1}{(1+\int_{\Omega}\frac{1}{1-\psi}\,\mathrm{d}x)^2}\frac{(2-u-\psi)v}{(1-u)^2(1-\psi)^2}\right].
\end{align*}
Since $\psi\in C^{\infty}(\Omega)$ is a steady-state solution and $\|u\|_{L^{\infty}}\le 1-\delta$, we have $\|\frac{1}{1-\psi}\|_{L^{\infty}}\le M$ and $\|\frac{1}{1-u}\|_{L^{\infty}}\le M$ for all $t\ge 0$.\\
\noindent
Thus
\begin{align*}
    &\quad\frac{1}{2}\frac{\mathrm{d}}{\mathrm{d}t}\|v\|_{L^2}^2+\|\nabla v\|_{L^2}^2\\
    &\ =\lambda\left(\int_{\Omega}\frac{v}{(1-u)^2}\frac{(2+\int_{\Omega}\frac{1}{1-u}\,\mathrm{d}x+\int_{\Omega}\frac{1}{1-\psi}\,\mathrm{d}x)(\int_{\Omega}\frac{-v}{(1-u)(1-\psi)}\,\mathrm{d}x)}{(1+\int_{\Omega}\frac{1}{1-u}\,\mathrm{d}x)^2(1+\int_{\Omega}\frac{1}{1-\psi}\,\mathrm{d}x)^2}\,\mathrm{d}x\right.\\
    &\quad\ \left.+\int_{\Omega}\frac{1}{(1+\int_{\Omega}\frac{1}{1-\psi}\,\mathrm{d}x)^2}\frac{(2-u-\psi)v^2}{(1-u)^2(1-\psi)^2}\,\mathrm{d}x\right)\\
    &\ \le C\left(\int_{\Omega}\frac{|v|}{(1-u)^2}\,\mathrm{d}x\int_{\Omega}\frac{|-v|}{(1-u)(1-\psi)}\,\mathrm{d}x+\int_{\Omega}\frac{(2-u-\psi)v^2}{(1-u)^2(1-\psi)^2}\,\mathrm{d}x\right)\\
    &\ \le C\|v\|_{L^2}^{2},
\end{align*}
yielding
\begin{equation}
    \frac{\mathrm{d}}{\mathrm{d}t}\|v\|_{L^2}^2+2\|\nabla v\|_{L^2}^2\le C\|v\|_{L^2}^{2}.\label{4.2}
\end{equation}
Multiplying (\ref{4.1}) by $v_t$ and integrating over $\Omega$:
\begin{align*}
    &\quad\frac{1}{2}\frac{\mathrm{d}}{\mathrm{d}t}\|\nabla v\|_{L^2}^2+\|v_t\|_{L^2}^2\\
    &\ =\lambda\left(\int_{\Omega}\frac{v_t}{(1-u)^2}\frac{(2+\int_{\Omega}\frac{1}{1-u}\,\mathrm{d}x+\int_{\Omega}\frac{1}{1-\psi}\,\mathrm{d}x)(\int_{\Omega}\frac{-v}{(1-u)(1-\psi)}\,\mathrm{d}x)}{(1+\int_{\Omega}\frac{1}{1-u}\,\mathrm{d}x)^2(1+\int_{\Omega}\frac{1}{1-\psi}\,\mathrm{d}x)^2}\,\mathrm{d}x\right.\\
    &\quad\ \left.+\int_{\Omega}\frac{1}{(1+\int_{\Omega}\frac{1}{1-\psi}\,\mathrm{d}x)^2}\frac{(2-u-\psi)v_t v}{(1-u)^2(1-\psi)^2}\,\mathrm{d}x\right)\\
    &\ \le C\left(\int_{\Omega}\frac{|v_t|}{(1-u)^2}\,\mathrm{d}x\int_{\Omega}\frac{|-v|}{(1-u)(1-\psi)}\,\mathrm{d}x+\int_{\Omega}\frac{(2-u-\psi)v_t v}{(1-u)^2(1-\psi)^2}\,\mathrm{d}x\right)\\
    &\ \le C\|v\|_{L^2}\|v_t\|_{L^2}\\
    &\ \le \frac{(C\|v\|_{L^2})^2}{2}+\frac{\|v_t\|_{L^2}^2}{2}.
\end{align*}
Yielding
\begin{equation}
    \frac{\mathrm{d}}{\mathrm{d}t}\|\nabla v\|_{L^2}^2+\|v_t\|_{L^2}^2\le C\|v\|_{L^2}^{2}.\label{4.3}
\end{equation}
Differentiating (\ref{4.1}) with respect to $t$:
\begin{align*}
  &\quad v_{tt}-\Delta v_t\\
  &\ =\frac{2\lambda u_t}{(1-u)^3(1+\int_{\Omega}\frac{1}{1-u}\,\mathrm{d}x)^2}-\frac{2\lambda\int_{\Omega}\frac{u_t}{(1-u)^2}\,\mathrm{d}x}{(1-u)^2(1+\int_{\Omega}\frac{1}{1-u}\,\mathrm{d}x)^3}-\frac{2\lambda \psi_t}{(1-\psi)^3(1+\int_{\Omega}\frac{1}{1-\psi}\,\mathrm{d}x)^2}\\
  &\quad\ +\frac{2\lambda\int_{\Omega}\frac{\psi_t}{(1-\psi)^2}\,\mathrm{d}x}{(1-\psi)^2(1+\int_{\Omega}\frac{1}{1-\psi}\,\mathrm{d}x)^3}.
\end{align*}
Since $\psi$ is a steady-state solution, $\psi_t=0$ and $v_t=u_t$, so
\begin{equation}
    v_{tt}-\Delta v_t=\frac{2\lambda v_t}{(1-u)^3(1+\int_{\Omega}\frac{1}{1-u}\,\mathrm{d}x)^2}-\frac{2\lambda\int_{\Omega}\frac{v_t}{(1-u)^2}\,\mathrm{d}x}{(1-u)^2(1+\int_{\Omega}\frac{1}{1-u}\,\mathrm{d}x)^3}.\label{4.4}
\end{equation}
Multiplying (\ref{4.4}) by $v_t$ and integrating over $\Omega$:
\begin{align*}
    &\quad \frac{1}{2}\frac{\mathrm{d}}{\mathrm{d}t}\|v_t\|_{L^2}^2+\|\nabla v_t\|_{L^2}^2\\
    &\ =\frac{2\lambda}{(1+\int_{\Omega}\frac{1}{1-u}\,\mathrm{d}x)^2}\int_{\Omega}\frac{v_{t}^{2}}{(1-u)^3}\,\mathrm{d}x-\frac{2\lambda(\int_{\Omega}\frac{v_t}{(1-u)^2}\,\mathrm{d}x)^2}{(1+\int_{\Omega}\frac{1}{1-u}\,\mathrm{d}x)^3}\\
    &\le \frac{2\lambda}{(1+\int_{\Omega}\frac{1}{1-u}\,\mathrm{d}x)^2}\int_{\Omega}\frac{v_{t}^{2}}{(1-u)^3}\,\mathrm{d}x\\
    &\le C_2\|v_t\|_{L^2}^{2},
\end{align*}
yielding
\begin{equation}
    \frac{1}{4C_2}\frac{\mathrm{d}}{\mathrm{d}t}\|v_t\|_{L^2}^2\le \frac{1}{2}\|v_t\|_{L^2}^{2}.\label{4.5}
\end{equation}
Adding (\ref{4.2}), (\ref{4.3}) and (\ref{4.5}):
$$\frac{\mathrm{d}}{\mathrm{d}t}\left(\|v\|_{L^2}^2+\|\nabla v\|_{L^2}^2+\frac{1}{4C_2}\|v_t\|_{L^2}^2\right)+2\|\nabla v\|_{L^2}^2+\frac{1}{2}\|v_t\|_{L^2}^2\le C\|v\|_{L^2}^{2},$$
thus
$$\frac{\mathrm{d}}{\mathrm{d}t}\left(\|v\|_{L^2}^2+\|\nabla v\|_{L^2}^2+\frac{1}{4C_2}\|v_t\|_{L^2}^2\right)+\|v\|_{L^2}^{2}+2\|\nabla v\|_{L^2}^2+\frac{1}{2}\|v_t\|_{L^2}^2\le (C+1)\|v\|_{L^2}^{2}.$$
Note there exists $\alpha>0$ such that
$$\|v\|_{L^2}^{2}+2\|\nabla v\|_{L^2}^2+\frac{1}{2}\|v_t\|_{L^2}^2\ge \alpha\left(\|v\|_{L^2}^2+\|\nabla v\|_{L^2}^2+\frac{1}{4C_2}\|v_t\|_{L^2}^2\right),$$
therefore
$$\frac{\mathrm{d}}{\mathrm{d}t}\widetilde{y}(t)+\alpha\widetilde{y}(t)\le C\|v\|_{L^2}^{2},$$
where $\widetilde{y}(t)=\|v\|_{L^2}^2+\|\nabla v\|_{L^2}^2+\frac{1}{4C_2}\|v_t\|_{L^2}^2$, yielding
\begin{equation}
    \frac{\,\mathrm{d}\widetilde{y}}{\,\mathrm{d}t}+\alpha\widetilde{y}\le C\|v\|_{L^2}^{2}.\label{4.6}
\end{equation}
When $\theta\in (0,\frac{1}{2})$, we have already shown $\|v\|_{L^2}\le C(1+t)^{-\frac{\theta}{1-2\theta}}$, so
$$\frac{\,\mathrm{d}\widetilde{y}}{\,\mathrm{d}t}+\alpha\widetilde{y}\le C(1+t)^{-\frac{2\theta}{1-2\theta}}.$$
Thus
\begin{align*}
    \widetilde{y}(t)&\le Ce^{-\alpha t}+Ce^{-\alpha t}\int_{0}^{t}(1+\tau)^{-\frac{2\theta}{1-2\theta}}e^{\alpha \tau}\,\mathrm{d}\tau\\
    &\le Ce^{-\alpha t}+Ce^{-\alpha t}\left[\int_{0}^{\frac{t}{2}}(1+\tau)^{-\frac{2\theta}{1-2\theta}}\,\mathrm{d}\tau\ e^{\frac{\alpha t}{2}}+(1+\frac{t}{2})^{-\frac{2\theta}{1-2\theta}}\int_{\frac{t}{2}}^{t}e^{\alpha\tau}\,\mathrm{d}\tau\right]\\
    &\le Ce^{-\alpha t}+Ce^{-\frac{\alpha t}{2}}\int_{0}^{\frac{t}{2}}(1+\tau)^{-\frac{2\theta}{1-2\theta}}\,\mathrm{d}\tau +Ce^{-\alpha t}(1+\frac{t}{2})^{-\frac{2\theta}{1-2\theta}}\left[\frac{1}{\alpha}(e^{\alpha t}-e^{\frac{\alpha t}{2}})\right]\\
    &\le C(1+t)^{-\frac{2\theta}{1-2\theta}}.
\end{align*}
Hence
$$\|\nabla v\|_{L^2}\le C(1+t)^{-\frac{\theta}{1-2\theta}},\quad \|v_t\|_{L^2}\le C(1+t)^{-\frac{\theta}{1-2\theta}},$$
and therefore
\begin{align*}
    \|v\|_{H^2}&\le C\|\Delta v\|_{L^2}\\
    &\le C\left(\|v_t\|_{L^2}+\left\|\frac{\lambda}{(1-u)^2(1+\int_{\Omega}\frac{1}{1-u}\,\mathrm{d}x)^2}-\frac{\lambda}{(1-\psi)^2(1+\int_{\Omega}\frac{1}{1-\psi}\,\mathrm{d}x)^2}\right\|_{L^2}\right)\\
    &\le C\|v_t\|_{L^2}+C\left\|\frac{1}{(1-u)^2}\int_{\Omega}\frac{-v}{(1-u)(1-\psi)}\,\mathrm{d}x\right\|_{L^2}+C\left\|\frac{(2-u-\psi)v}{(1-u)^2(1-\psi)^2}\right\|_{L^2}\\
    &\le C\|v_t\|_{L^2}+C\|v\|_{L^2}\\
    &\le C(1+t)^{-\frac{\theta}{1-2\theta}}.
\end{align*}
When $\theta=\frac{1}{2}$, we have $\|v\|_{L^2}\le C_0 e^{-C_1 t}$, so
$$\frac{\,\mathrm{d}\widetilde{y}}{\,\mathrm{d}t}+\alpha\widetilde{y}\le Ce^{-2C_1 t}.$$
Thus
\begin{align*}
    \widetilde{y}(t)&\le Ce^{-\alpha t}+Ce^{-\alpha t}\int_{0}^{t}e^{(\alpha-2C_1)\tau}\,\mathrm{d}\tau\\
    &= Ce^{-\alpha t}+Ce^{-\alpha t}\frac{1}{\alpha-2C_1}\left(e^{(\alpha-2C_1)t}-1\right)\\
    &\le Ce^{-\alpha t}+Ce^{-2C_1 t}+Ce^{-\alpha t}\\
    &\le Ce^{-\min\{\alpha,2C_1\}t}.
\end{align*}
Hence
$$\|\nabla v\|_{L^2}\le Ce^{-\min\{\frac{\alpha}{2},C_1\}t},\quad \|v_t\|_{L^2}\le Ce^{-\min\{\frac{\alpha}{2},C_1\}t},$$
and similarly
$$\|v\|_{H^2}\le C\|\Delta v\|_{L^2}\le C\|v_t\|_{L^2}+C\|v\|_{L^2}\le Ce^{-\min\{\frac{\alpha}{2},C_1\}t}.$$
The theorem is proved.
\end{proof}

\end{document}
